% ============================================================================
% 30_ai.tex  --  AI-collaboration note  (~1/3 page)
%
% GOAL: A short, honest reflection on how LLMs were used as a research
%       collaborator for this part (what they helped derive/debug/write, and
%       where you corrected them). This is a course requirement of the
%       AI-augmented seminar.
% ============================================================================

\section{On the AI collaboration}
\label{sec:ai}

My experience using large language models throughout the project was generally
positive.  We used Anthropic's Claude Opus 4.8 through Claude Code and OpenAI's
GPT-5.6 Sol through Codex CLI.  The models helped us to prepare the weekly presentations,
write topic-specific notes, and solve the main coding challenges.  They proved to be helpful as
editors and as personal physics teachers, especially when generating all the learning materials.  In addition, a collaborative GitHub repository and regular
group meetings also contributed to a professional working environment.  The
project inspired me to learn more about quantum computing and recent
developments in topological quantum computers.
